\documentclass[11pt, a4paper]{article}

\usepackage[a4paper, margin=2.5cm]{geometry}
\usepackage[brazil]{babel}
\usepackage[utf8]{inputenc}
\usepackage[T1]{fontenc}

\usepackage{amsmath}
\usepackage{booktabs}
\usepackage{enumitem}
\setlist[itemize]{label=-}
\usepackage{xcolor}
\usepackage{listings}
\usepackage{hyperref}

\hypersetup{
    colorlinks=true,
    linkcolor=blue!70!black,
    citecolor=blue!70!black,
    urlcolor=blue!70!black
}

\lstset{
    basicstyle=\ttfamily\small,
    breaklines=true,
    frame=single,
    backgroundcolor=\color{gray!10},
    numbers=none
}

\title{Atividade 1 - Mineração de Dados}
\author{Gabriel M. Brum,  Pedro H. M. Gatti  \\ Programa da Pós-Graduação em Ciência da Computação (PPGCC) - UNESP}
\date{\today}

\begin{document}

\maketitle

\begin{abstract}
resumo
\end{abstract}

\vspace{0.5cm}

\section{Introdução}

Nesta atividade foi solicitado escolhermos um conjunto de dados que contenha ao menos 4 classes, então aplicar o pré-processamento que julgassemos adequado. Após isso, deveríamos testar os seguintes algoritmos: árvore, \textit{bagging}, \textit{boosting} e \textit{random forest}. Ao final, avaliar quais deles gerou o melhor modelo usando tanto alguma técnica de validação, escolhendo uma métrica de desempenho. Considerando o de melhor resultado, faça a \textit{setagem} de hiperparâmetros no mesmo para ver se o desempenho melhora. Adicionalmente, deveriamos transformar o problema multiclasse em problemas binários usando OVA e OVO. Escolher uma métrica de desempenho e analisar qual das abordagens se saiu melhor para alguma técnica de avaliação utilizando o algoritmo de árvore.

\section{\textit{Dataset}}

O conjunto de dados \textit{Car Evaluation} é um \textit{benchmark} da área de aprendizado de máquina, disponibilizado no repositório UCI por Bohanec e Zupan em 1997 \cite{zupan1997machine}. Sua origem remonta aos trabalhos de Bohanec e Rajkovič (1988, 1990) no contexto de sistemas especialistas de apoio à decisão multicritério (ferramenta DEX) \cite{bohanec1988knowledge}. Originalmente, o domínio foi modelado através de uma estrutura hierárquica na qual seis atributos primários determinavam conceitos intermediários de custo e viabilidade técnica antes de inferir a aceitabilidade final do veículo. Na versão publicada, os conceitos intermediários foram suprimidos, restando uma relação direta entre os seis atributos descritivos e a classe-alvo distribuída em 1728 instâncias (cobrindo a totalidade do espaço vetorial discreto do problema).

O conjunto tem um total de 1728 instâncias, 6 atributos e 4 classes (\textit{unacc = unnacceptable}, \textit{acc = acceptable}, \textit{good} e \textit{vgood = very good}) que indicam a condição de compra de veículos. A distribuição de classes está apresentada na Tabela \ref{tab:distribuicao_classes}, já os atributos estão descritos na Tabela \ref{tab:descricao_atributos}.

\begin{table}[htbp]
    \centering
    \caption{Distribuição das classes no conjunto de dados Car Evaluation.}
    \label{tab:distribuicao_classes}
    \begin{tabular}{|l|c|r|}
        \hline
        Classe & Quantidade de instâncias & \% do total de amostras \\
        \hline
        \textit{unacc}  & 1210 & 70{,}023\% \\
        \textit{acc}    & 384  & 22{,}222\% \\
        \textit{good}   & 69   & 3{,}993\%  \\
        \textit{vgood}  & 65   & 3{,}762\%  \\
        \hline
        \textbf{Total}  & 1728 & 100{,}000\% \\
        \hline
    \end{tabular}
\end{table}

\begin{table}[htbp]
    \centering
    \caption{Descrição dos atributos e variável-alvo do conjunto de dados Car Evaluation.}
    \label{tab:descricao_atributos}
    \begin{tabular}{llll}
        \hline
        \textbf{Atributo} &  \textbf{Valores} & \textbf{Ordem Natural} \\
        \hline
        \texttt{buying}   &  \textit{low}, \textit{med}, \textit{high}, \textit{v-high} & $\text{low} < \text{med} < \text{high} < \text{v-high}$ \\
        \texttt{maint}    &  \textit{low}, \textit{med}, \textit{high}, \textit{v-high} & $\text{low} < \text{med} < \text{high} < \text{v-high}$ \\
        \texttt{doors}    &  \textit{2}, \textit{3}, \textit{4}, \textit{5more} & $2 < 3 < 4 < \text{5more}$ \\
        \texttt{persons}  &  \textit{2}, \textit{4}, \textit{more} & $2 < 4 < \text{more}$ \\
        \texttt{lug\_boot} &  \textit{small}, \textit{med}, \textit{big} & $\text{small} < \text{med} < \text{big}$ \\
        \texttt{safety}   &  \textit{low}, \textit{med}, \textit{high} & $\text{low} < \text{med} < \text{high}$ \\
        \hline
        \textbf{Classe} (\texttt{class}) &  \textit{unacc}, \textit{acc}, \textit{good}, \textit{vgood} & $\text{unacc} < \text{acc} < \text{good} < \text{vgood}$ \\
        \hline
    \end{tabular}
\end{table}

\section{Pré-Processamento}

Por conta de se tratar de um conjunto de dados categórico, e que possuia ordem entre as classes, optamos por utilizar a técnica de \textit{Ordinal Encoding} para transformar os atributos em valores numéricos. A Tabela \ref{tab:label_encoding} apresenta a codificação utilizada para cada atributo.

\begin{table}[htbp]
    \centering
    \caption{Mapeamento de codificação ordinal aplicado aos atributos e à classe-alvo.}
    \label{tab:label_encoding}
    \begin{tabular}{lll}
        \hline
        \textbf{Atributo} & \textbf{Categoria Original} & \textbf{Valor Codificado} \\
        \hline
        \texttt{buying}   & \textit{low}, \textit{med}, \textit{high}, \textit{vhigh}   & 1, 2, 3, 4 \\
        \texttt{maint}    & \textit{low}, \textit{med}, \textit{high}, \textit{vhigh}   & 1, 2, 3, 4 \\
        \texttt{doors}    & \textit{2}, \textit{3}, \textit{4}, \textit{5more}          & 2, 3, 4, 5 \\
        \texttt{persons}  & \textit{2}, \textit{4}, \textit{more}                       & 2, 4, 5 \\
        \texttt{lug\_boot} & \textit{small}, \textit{med}, \textit{big}                 & 1, 2, 3 \\
        \texttt{safety}   & \textit{low}, \textit{med}, \textit{high}                   & 1, 2, 3 \\
        \hline
        \textbf{Classe} (\texttt{class}) & \textit{unacc}, \textit{acc}, \textit{good}, \textit{vgood} & 0, 1, 2, 3 \\
        \hline
    \end{tabular}
\end{table}

\section{Metodologia}
\label{sec:metodologia}

O procedimento experimental foi estruturado em duas etapas principais, orientadas pelo uso de validação estratificada e pela métrica \textit{Macro F1-Score} ($F1_{\text{macro}}$), fundamental para lidar com o desbalanceamento das quatro classes do problema.

\subsection{Parte I: Comparação de Modelos e Otimização de Hiperparâmetros}
\label{subsec:metodologia_parte1}

Nesta etapa, comparou-se o desempenho preditivo de quatro algoritmos de aprendizado de máquina supervisionado:
\begin{itemize}
    \item \textbf{Árvore de Decisão (\textit{Decision Tree}):} Classificador base baseado em divisões recursivas de atributos;
    \item \textbf{Bagging (\textit{Bootstrap Aggregating}):} Comitê que combina estimadores independentes treinados com amostragem com reposição;
    \item \textbf{AdaBoost (\textit{Adaptive Boosting}):} Método de agregação sequencial que ajusta pesos iterativamente para instâncias classificadas incorretamente;
    \item \textbf{Random Forest:} Floresta aleatória que combina amostragem \textit{bootstrap} com seleção aleatória de subconjuntos de atributos em cada nó.
\end{itemize}

O protocolo de validação adotado para a comparação inicial dos quatro modelos foi a validação cruzada estratificada em 5 \textit{folds} (\textit{Stratified 5-Fold Cross-Validation}) com embaralhamento dos dados (\textit{shuffle}), utilizando o $F1_{\text{macro}}$ como critério decisor.

O classificador que alcançou o maior valor médio de $F1_{\text{macro}}$ foi selecionado para a fase de ajuste fino de hiperparâmetros (\textit{hyperparameter tuning}). O ajuste foi conduzido através de \textit{GridSearchCV} com validação cruzada interna estratificada de 5 \textit{folds} e métrica de pontuação \texttt{scoring='f1\_macro'}. Por fim, avaliou-se o ganho de desempenho obtido confrontando os resultados do modelo antes e após a otimização no conjunto de teste independente (Holdout 70/30 estratificado).

\subsection{Parte II: Decomposição Multiclasse em Problemas Binários}
\label{subsec:metodologia_parte2}

Nesta etapa, investigou-se a estratégia de binarização do problema multiclasse utilizando exclusivamente a Árvore de Decisão (\texttt{DecisionTreeClassifier(random\_state=42)}) como algoritmo base:

\begin{itemize}
    \item \textbf{One-vs-All (OVA / One-vs-Rest):} Implementado por meio da classe \texttt{OneVsRestClassifier}, treina um classificador binário para cada classe contra o agrupamento das demais, totalizando 4 modelos;
    \item \textbf{One-vs-One (OVO):} Implementado por meio da classe \texttt{OneVsOneClassifier}, treina um classificador binário dedicado para cada par possível de classes, totalizando $\frac{4 \times 3}{2} = 6$ modelos, com resolução final por votação majoritária.
\end{itemize}

Ambas as abordagens foram avaliadas sob dois esquemas de validação complementares:
\begin{enumerate}
    \item \textbf{Holdout Estratificado (70/30):} Particionamento dos dados em 70\% para treino e 30\% para teste (\texttt{random\_state=42}), preservando a proporção de cada classe. A partir das predições no conjunto de teste, computaram-se a acurácia, o $F1_{\text{macro}}$ e as matrizes de confusão de cada método.
    \item \textbf{Validação Cruzada Estratificada ($k=5$):} Avaliação via \texttt{StratifiedKFold(n\_splits=5, shuffle=True, random\_state=42)}, mensurando a média e o desvio padrão da acurácia e do $F1_{\text{macro}}$ ao longo das 5 dobras para determinar a abordagem de melhor desempenho.
\end{enumerate}

\bibliographystyle{plain} % ou 'abntex2-alf', 'ieee', etc.
\bibliography{referencias} % Nome do seu arquivo referencias.bib (sem a extensão)


\end{document}